\documentclass[12pt,a4paper]{exam}
\usepackage[utf8]{inputenc}
\usepackage{amsmath,amssymb,amsthm}
\usepackage[margin=1in]{geometry}
\usepackage{graphicx}
\usepackage[colorlinks=true]{hyperref}
\pagestyle{headandfoot}
\firstpageheader{MIT OpenCourseWare}{}{\textbf{EXTREMELY HARD -- MIT LEVEL}}
\runningheader{MIT OpenCourseWare}{}{\textbf{EXTREMELY HARD -- MIT LEVEL}}
\firstpagefooter{}{Page \thepage}{}
\runningfooter{}{Page \thepage}{}

\begin{document}

\begin{center}
{\LARGE \textbf{MIT OpenCourseWare}}\\8.321\\Quantum Field Theory I
\vspace{0.5cm}
{8.321} --- {Quantum Field Theory I}
\vspace{0.3cm}
May 2026
\vspace{0.3cm}
\textbf{EXTREMELY HARD -- MIT LEVEL}
\end{center}

\begin{center}
\textbf{Time Limit: 3 hours}
\end{center}

\begin{center}
\textbf{Total Marks: 100}
\end{center}

\vspace{0.5cm}

\begin{questions}

\question[4]
Construct an original proof-level problem involving Canonical Quantization for 8.321 (Quantum Field Theory I) and provide a complete solution. Your problem should be at the level of a graduate qualifying exam.

\vspace{0.3cm}

\question[5]
Construct an original proof-level problem involving Path Integrals for 8.321 (Quantum Field Theory I) and provide a complete solution. Your problem should be at the level of a graduate qualifying exam.

\vspace{0.3cm}

\question[5]
Prove the fundamental theorem concerning Gauge Fields in the context of 8.321 (Quantum Field Theory I). State the theorem precisely, provide a complete proof, and discuss three important corollaries.

\vspace{0.3cm}

\question[4]
Prove the fundamental theorem concerning Spontaneous Symmetry Breaking in the context of 8.321 (Quantum Field Theory I). State the theorem precisely, provide a complete proof, and discuss three important corollaries.

\vspace{0.3cm}

\question[4]
Prove the fundamental theorem concerning Renormalization in the context of 8.321 (Quantum Field Theory I). State the theorem precisely, provide a complete proof, and discuss three important corollaries.

\vspace{0.3cm}

\question[6]
Analyze the limitations of standard approaches to Canonical Quantization when applied to edge cases in 8.321. Construct explicit counterexamples showing where intuitive reasoning fails.

\vspace{0.3cm}

\question[5]
Construct an original proof-level problem involving Path Integrals for 8.321 (Quantum Field Theory I) and provide a complete solution. Your problem should be at the level of a graduate qualifying exam.

\vspace{0.3cm}

\question[2]
Construct an original proof-level problem involving Gauge Fields for 8.321 (Quantum Field Theory I) and provide a complete solution. Your problem should be at the level of a graduate qualifying exam.

\vspace{0.3cm}

\question[2]
Analyze the limitations of standard approaches to Spontaneous Symmetry Breaking when applied to edge cases in 8.321. Construct explicit counterexamples showing where intuitive reasoning fails.

\vspace{0.3cm}

\question[2]
Analyze the limitations of standard approaches to Renormalization when applied to edge cases in 8.321. Construct explicit counterexamples showing where intuitive reasoning fails.

\vspace{0.3cm}

\question[6]
Prove the fundamental theorem concerning Canonical Quantization in the context of 8.321 (Quantum Field Theory I). State the theorem precisely, provide a complete proof, and discuss three important corollaries.

\vspace{0.3cm}

\question[3]
Prove the fundamental theorem concerning Path Integrals in the context of 8.321 (Quantum Field Theory I). State the theorem precisely, provide a complete proof, and discuss three important corollaries.

\vspace{0.3cm}

\question[5]
Construct an original proof-level problem involving Gauge Fields for 8.321 (Quantum Field Theory I) and provide a complete solution. Your problem should be at the level of a graduate qualifying exam.

\vspace{0.3cm}

\question[2]
Prove the fundamental theorem concerning Spontaneous Symmetry Breaking in the context of 8.321 (Quantum Field Theory I). State the theorem precisely, provide a complete proof, and discuss three important corollaries.

\vspace{0.3cm}

\question[3]
Construct an original proof-level problem involving Renormalization for 8.321 (Quantum Field Theory I) and provide a complete solution. Your problem should be at the level of a graduate qualifying exam.

\vspace{0.3cm}

\question[4]
Analyze the limitations of standard approaches to Canonical Quantization when applied to edge cases in 8.321. Construct explicit counterexamples showing where intuitive reasoning fails.

\vspace{0.3cm}

\question[4]
Analyze the limitations of standard approaches to Path Integrals when applied to edge cases in 8.321. Construct explicit counterexamples showing where intuitive reasoning fails.

\vspace{0.3cm}

\question[5]
Analyze the limitations of standard approaches to Gauge Fields when applied to edge cases in 8.321. Construct explicit counterexamples showing where intuitive reasoning fails.

\vspace{0.3cm}

\question[3]
Construct an original proof-level problem involving Spontaneous Symmetry Breaking for 8.321 (Quantum Field Theory I) and provide a complete solution. Your problem should be at the level of a graduate qualifying exam.

\vspace{0.3cm}

\question[6]
Analyze the limitations of standard approaches to Renormalization when applied to edge cases in 8.321. Construct explicit counterexamples showing where intuitive reasoning fails.

\vspace{0.3cm}

\question[4]
Prove the fundamental theorem concerning Canonical Quantization in the context of 8.321 (Quantum Field Theory I). State the theorem precisely, provide a complete proof, and discuss three important corollaries.

\vspace{0.3cm}

\question[3]
Analyze the limitations of standard approaches to Path Integrals when applied to edge cases in 8.321. Construct explicit counterexamples showing where intuitive reasoning fails.

\vspace{0.3cm}

\question[6]
Construct an original proof-level problem involving Gauge Fields for 8.321 (Quantum Field Theory I) and provide a complete solution. Your problem should be at the level of a graduate qualifying exam.

\vspace{0.3cm}

\question[3]
Prove the fundamental theorem concerning Spontaneous Symmetry Breaking in the context of 8.321 (Quantum Field Theory I). State the theorem precisely, provide a complete proof, and discuss three important corollaries.

\vspace{0.3cm}

\question[4]
Analyze the limitations of standard approaches to Renormalization when applied to edge cases in 8.321. Construct explicit counterexamples showing where intuitive reasoning fails.

\vspace{0.3cm}

\end{questions}

\newpage
\section*{Solutions}

\textbf{Solution 1:}

Solution approach: First recall the definitions and key theorems related to Canonical Quantization from 8.321. Then apply the appropriate techniques to construct a rigorous argument. The solution must be self-contained and reference only the material from the course.

\vspace{0.5cm}

\textbf{Solution 2:}

Solution approach: First recall the definitions and key theorems related to Path Integrals from 8.321. Then apply the appropriate techniques to construct a rigorous argument. The solution must be self-contained and reference only the material from the course.

\vspace{0.5cm}

\textbf{Solution 3:}

The proof follows from the definitions and properties established in 8.321. The key insight is to apply the relevant theorem and verify the hypotheses hold. Detailed solution depends on the specific formulation of the theorem.

\vspace{0.5cm}

\textbf{Solution 4:}

The proof follows from the definitions and properties established in 8.321. The key insight is to apply the relevant theorem and verify the hypotheses hold. Detailed solution depends on the specific formulation of the theorem.

\vspace{0.5cm}

\textbf{Solution 5:}

The proof follows from the definitions and properties established in 8.321. The key insight is to apply the relevant theorem and verify the hypotheses hold. Detailed solution depends on the specific formulation of the theorem.

\vspace{0.5cm}

\textbf{Solution 6:}

The edge case analysis reveals the subtle assumptions underlying standard treatments of Canonical Quantization. By constructing explicit counterexamples, we see that the usual hypotheses cannot be weakened without loss of the theorem's conclusion.

\vspace{0.5cm}

\textbf{Solution 7:}

Solution approach: First recall the definitions and key theorems related to Path Integrals from 8.321. Then apply the appropriate techniques to construct a rigorous argument. The solution must be self-contained and reference only the material from the course.

\vspace{0.5cm}

\textbf{Solution 8:}

Solution approach: First recall the definitions and key theorems related to Gauge Fields from 8.321. Then apply the appropriate techniques to construct a rigorous argument. The solution must be self-contained and reference only the material from the course.

\vspace{0.5cm}

\textbf{Solution 9:}

The edge case analysis reveals the subtle assumptions underlying standard treatments of Spontaneous Symmetry Breaking. By constructing explicit counterexamples, we see that the usual hypotheses cannot be weakened without loss of the theorem's conclusion.

\vspace{0.5cm}

\textbf{Solution 10:}

The edge case analysis reveals the subtle assumptions underlying standard treatments of Renormalization. By constructing explicit counterexamples, we see that the usual hypotheses cannot be weakened without loss of the theorem's conclusion.

\vspace{0.5cm}

\textbf{Solution 11:}

The proof follows from the definitions and properties established in 8.321. The key insight is to apply the relevant theorem and verify the hypotheses hold. Detailed solution depends on the specific formulation of the theorem.

\vspace{0.5cm}

\textbf{Solution 12:}

The proof follows from the definitions and properties established in 8.321. The key insight is to apply the relevant theorem and verify the hypotheses hold. Detailed solution depends on the specific formulation of the theorem.

\vspace{0.5cm}

\textbf{Solution 13:}

Solution approach: First recall the definitions and key theorems related to Gauge Fields from 8.321. Then apply the appropriate techniques to construct a rigorous argument. The solution must be self-contained and reference only the material from the course.

\vspace{0.5cm}

\textbf{Solution 14:}

The proof follows from the definitions and properties established in 8.321. The key insight is to apply the relevant theorem and verify the hypotheses hold. Detailed solution depends on the specific formulation of the theorem.

\vspace{0.5cm}

\textbf{Solution 15:}

Solution approach: First recall the definitions and key theorems related to Renormalization from 8.321. Then apply the appropriate techniques to construct a rigorous argument. The solution must be self-contained and reference only the material from the course.

\vspace{0.5cm}

\textbf{Solution 16:}

The edge case analysis reveals the subtle assumptions underlying standard treatments of Canonical Quantization. By constructing explicit counterexamples, we see that the usual hypotheses cannot be weakened without loss of the theorem's conclusion.

\vspace{0.5cm}

\textbf{Solution 17:}

The edge case analysis reveals the subtle assumptions underlying standard treatments of Path Integrals. By constructing explicit counterexamples, we see that the usual hypotheses cannot be weakened without loss of the theorem's conclusion.

\vspace{0.5cm}

\textbf{Solution 18:}

The edge case analysis reveals the subtle assumptions underlying standard treatments of Gauge Fields. By constructing explicit counterexamples, we see that the usual hypotheses cannot be weakened without loss of the theorem's conclusion.

\vspace{0.5cm}

\textbf{Solution 19:}

Solution approach: First recall the definitions and key theorems related to Spontaneous Symmetry Breaking from 8.321. Then apply the appropriate techniques to construct a rigorous argument. The solution must be self-contained and reference only the material from the course.

\vspace{0.5cm}

\textbf{Solution 20:}

The edge case analysis reveals the subtle assumptions underlying standard treatments of Renormalization. By constructing explicit counterexamples, we see that the usual hypotheses cannot be weakened without loss of the theorem's conclusion.

\vspace{0.5cm}

\textbf{Solution 21:}

The proof follows from the definitions and properties established in 8.321. The key insight is to apply the relevant theorem and verify the hypotheses hold. Detailed solution depends on the specific formulation of the theorem.

\vspace{0.5cm}

\textbf{Solution 22:}

The edge case analysis reveals the subtle assumptions underlying standard treatments of Path Integrals. By constructing explicit counterexamples, we see that the usual hypotheses cannot be weakened without loss of the theorem's conclusion.

\vspace{0.5cm}

\textbf{Solution 23:}

Solution approach: First recall the definitions and key theorems related to Gauge Fields from 8.321. Then apply the appropriate techniques to construct a rigorous argument. The solution must be self-contained and reference only the material from the course.

\vspace{0.5cm}

\textbf{Solution 24:}

The proof follows from the definitions and properties established in 8.321. The key insight is to apply the relevant theorem and verify the hypotheses hold. Detailed solution depends on the specific formulation of the theorem.

\vspace{0.5cm}

\textbf{Solution 25:}

The edge case analysis reveals the subtle assumptions underlying standard treatments of Renormalization. By constructing explicit counterexamples, we see that the usual hypotheses cannot be weakened without loss of the theorem's conclusion.

\vspace{0.5cm}

\end{document}